%==============================================================================
\section{Xây dựng đặc trưng tiến (Forward Feature Construction)}
\label{sec:forward-feature-construction}
%==============================================================================

\begin{dinhnghia}[title={Forward Feature Construction}]
Xây dựng đặc trưng tiến là phương pháp \textbf{lựa chọn đặc trưng}: bắt đầu từ tập đặc trưng \textbf{rỗng}, rồi \textbf{lặp thêm} đặc trưng cho hiệu năng tốt nhất ở mỗi bước cho đến khi đạt số đặc trưng mong muốn.
\end{dinhnghia}

\begin{ynghia}[title={target lựa chọn đặc trưng}]
target là xác định các đặc trưng \textbf{quan trọng nhất} để dự đoán biến target, đồng thời bỏ qua đặc trưng ít quan trọng --- vốn thêm \textbf{nhiễu} vào mô hình và có thể gây \textbf{overfitting}.
\end{ynghia}

\subsection{Các bước thuật toán}

\begin{phuongphap}[title={Các bước Forward Feature Construction}]
\begin{enumerate}
  \item \textbf{Khởi tạo} tập đặc trưng rỗng.
  \item \textbf{Đặt} số đặc trưng tối đa cần chọn.
  \item \textbf{Lặp} cho đến khi đủ số đặc trưng:
  \begin{itemize}
    \item Với mỗi đặc trưng \textbf{chưa} nằm trong tập đã chọn: khớp mô hình gồm tập đã chọn \textbf{cộng} đặc trưng hiện tại, đánh giá hiệu năng trên tập \textbf{kiểm tra} (validation/test).
    \item \textbf{Chọn} đặc trưng cho hiệu năng tốt nhất và \textbf{thêm} vào tập đã chọn.
  \end{itemize}
  \item \textbf{Trả về} tập đặc trưng đã chọn làm tập tối ưu cho mô hình.
\end{enumerate}
\end{phuongphap}

\begin{tinhchat}[title={Ưu điểm}]
Ưu điểm chính: \textbf{hiệu quả tính toán}, có thể áp dụng cho dữ liệu \textbf{chiều cao}.
\end{tinhchat}

\begin{luuy}[title={Hạn chế}]
Không luôn dẫn tới tập đặc trưng \textbf{tối ưu toàn cục}, đặc biệt khi có đặc trưng \textbf{tương quan cao} hoặc quan hệ \textbf{phi tuyến} giữa đặc trưng và biến target.
\end{luuy}

\subsection{Ví dụ triển khai Python}

\begin{dongco}[title={Bối cảnh}]
Ví dụ dưới dùng hồi quy tuyến tính (\texttt{LinearRegression}), điểm \texttt{score} trên tập kiểm tra làm tiêu chí chọn đặc trưng tiến.
\end{dongco}

\begin{vidu}[title={Code Python đầy đủ}]
\begin{lstlisting}
# Import thu vien
import pandas as pd
import numpy as np
from sklearn.model_selection import train_test_split
from sklearn.linear_model import LinearRegression
from sklearn.datasets import fetch_openml

# Tai bo diabetes
diabetes = fetch_openml('diabetes', version=0, as_frame=True, parser='auto')
X = diabetes.iloc[:, :-1].values
y = diabetes.iloc[:, -1].astype(int).values

# Chia tap huan luyen / kiem tra
X_train, X_test, y_train, y_test = train_test_split(
    X, y, test_size=0.2, random_state=0)

# Tap dac trung da chon (rong)
selected_features = set()

# So dac trung toi da
max_features = 8

# Lap den khi du 8 dac trung
while len(selected_features) < max_features:

    best_feature = None
    best_score = 0

    for i in range(X_train.shape[1]):

        if i in selected_features:
            continue

        X_train_selected = X_train[:, list(selected_features) + [i]]
        regressor = LinearRegression()
        regressor.fit(X_train_selected, y_train)

        X_test_selected = X_test[:, list(selected_features) + [i]]
        score = regressor.score(X_test_selected, y_test)

        if score > best_score:
            best_feature = i
            best_score = score

    selected_features.add(best_feature)

    print('Selected Features:', list(selected_features))
    print('Score:', best_score)
\end{lstlisting}
\end{vidu}

\begin{ynghia}[title={Cách đọc vòng lặp}]
Ở mỗi vòng \texttt{while}, thuật toán thử \textbf{từng} cột chưa chọn: huấn luyện trên \texttt{X\_train} gồm các cột đã chọn + cột thử \texttt{i}, tính \texttt{regressor.score} trên \texttt{X\_test} tương ứng. Cột có \texttt{score} cao nhất được \textbf{cố định} vào \texttt{selected\_features}. In ra tập chỉ số và điểm sau mỗi bước.
\end{ynghia}

\begin{vidu}[title={Kết quả chạy (mẫu)}]
Khi chạy, terminal có thể in:
\begin{lstlisting}
Selected Features: [1]
Score: 0.23530716168783583
Selected Features: [0, 1]
Score: 0.2923143573608237
Selected Features: [0, 1, 5]
Score: 0.3164103491569179
Selected Features: [0, 1, 5, 6]
Score: 0.3287368302427327
Selected Features: [0, 1, 2, 5, 6]
Score: 0.334586804842275
Selected Features: [0, 1, 2, 3, 5, 6]
Score: 0.3356264736550455
Selected Features: [0, 1, 2, 3, 4, 5, 6]
Score: 0.3313166516703744
Selected Features: [0, 1, 2, 3, 4, 5, 6, 7]
Score: 0.32230203252064216
\end{lstlisting}
\end{vidu}

\begin{luuy}[title={Chỉ số đặc trưng}]
Các số \texttt{0, 1, 2, \ldots} là \textbf{chỉ số cột} trong ma trận \texttt{X} (thứ tự cột của bộ Pima sau \texttt{fetch\_openml}), không phải tên biến. Thứ tự thêm phụ thuộc dữ liệu và \texttt{random\_state} khi chia train/test.
\end{luuy}
